\section{アプローチ01（ベースライン手法）}

\subsection{手法の概要}
アプローチ01は、本研究における比較の基準となるベースライン手法である。監視カメラ等から取得された標準的な二次元画像フレーム全体をそのまま入力とし、一般的なオブジェクト検出モデルを用いて画面内に存在する人物（\texttt{person} クラス）を一元的に列挙・検出する。

得られた二次元バウンディングボックスの中心座標間におけるユーグリッド距離を測定し、事前に設定した一定の物理距離閾値（$\delta_{dist}$）以下に位置する人物同士を同一のグループとして簡易的に統合・判定する手法を採用した。

\subsection{結果と課題}
ベースライン手法の検証結果および明瞭となった課題は以下の通りである。

\begin{enumerate}
    \item \textbf{画角全体のノイズと通行人の誤検出}：
    待機列領域以外の画面周辺部（背景、通路、他展示のスペース）に存在する無関係な人物や通行人もすべて一律に検出対象となるため、待機列人数の正確な計測を阻害する大きなノイズとなった。
    \item \textbf{重なり・遠近感による距離判定の破綻}：
    カメラ画像上の二次元距離のみを基準とするため、奥行き方向（カメラからの距離）の違いによるスケール変化や、待機列の折返しによって物理的には離れている前後の人物が画面上では近接して配置される場合、誤って同一グループとして誤統合される事態が頻発した。
    \item \textbf{遮蔽に対する脆弱性}：
    密集部での人物同士の重なり（オクルージョン）によりバウンディングボックスの検出自体がスキップされるか、あるいは極端に大きなバウンディングボックスが1つだけ生成され、個体認識が機能しない課題が判明した。
\end{enumerate}
